\section{Throughput Analysis by Corridor}
\label{sec:throughput}

\subsection{Current V/C Ratios from HPMS Data}

Table~\ref{tab:vc-current} presents current peak-period V/C ratios for all 8 T1 corridors, derived from the Highway Performance Monitoring System (HPMS) data set \citep{FHWA_HPMS2023}. V/C ratios are computed at the 90th-percentile segment (the segment with the highest V/C in each corridor, representing the binding constraint on freight throughput). All values represent weekday peak-period (6 AM--10 PM) directional V/C.

\begin{table}[h!]
\centering
\caption{T1 Corridor Peak V/C Ratios and Managed Lane Program Status}
\label{tab:vc-current}
\begin{tabular}{llrrrl}
\toprule
\textbf{Corridor} & \textbf{Location (90th pct segment)} & \textbf{V/C} & \textbf{PTI} & \textbf{AADT (M)} & \textbf{ML status} \\
\midrule
I-90  & Chicago to Milwaukee (Tri-State)    & 2.21 & 2.18 & 187 & Priority \\
I-80  & Bay Area (I-580 to I-880 merge)     & 2.23 & 2.09 & 172 & Priority \\
I-10  & Houston (I-610 to I-45 S)           & 2.11 & 2.02 & 196 & Priority \\
I-95  & New Jersey (Exit 12 to Exit 8A)     & 1.98 & 1.94 & 214 & Priority \\
I-75  & Atlanta (I-285 S to I-285 N)        & 1.87 & 1.83 & 178 & Priority \\
I-85  & Charlotte (I-485 N to US-74)        & 1.74 & 1.71 & 143 & Priority \\
I-35  & Dallas--Fort Worth (I-635 to US-380) & 1.62 & 1.58 & 168 & Priority \\
I-40  & Statewide maximum (Oklahoma City)   & 0.84 & 1.11 &  89 & Excluded \\
\bottomrule
\end{tabular}
\end{table}

All 7 priority corridors exceed V/C = 1.0 at the 90th percentile segment --- meaning they are in chronic LOS F (oversaturated) during peak periods. The PTI values of 1.58--2.21 indicate that shippers must pad transit time by 58--121\% above free-flow to guarantee on-time delivery at the 95th percentile. At a cross-country freight rate of \$3.50--\$4.50 per mile, this transit time uncertainty imposes a hedging cost equivalent to 12--18\% of freight revenue \citep{ATRI_costs2024}.

\subsection{I-40: The Exception That Validates the Rule}

Interstate 40 is the one T1 corridor excluded from the managed freight lane program. Its peak V/C of 0.84 at the statewide maximum (Oklahoma City) is at or below the I2.0 capacity target of V/C $\leq$ 0.85. The corridor operates at LOS C--D across its full length; its PTI of 1.11 indicates that shippers need buffer only 11\% above free-flow transit time for on-time delivery --- well within acceptable commercial tolerances.

The I-40 exclusion is significant for two reasons. First, it provides external validation of the ROUTE tier methodology. Paper B.2 \citep{ROUTE_B2} found that I-40 has zero ATRI top-50 bottleneck locations --- an independent confirmation that the corridor is not under capacity stress. The HPMS V/C analysis and the ATRI revealed-preference data converge on the same finding from independent methodologies.

Second, the I-40 exclusion demonstrates that the I2.0 managed lane program is demand-driven, not tier-driven. Tier classification determines strategic importance; capacity analysis determines whether managed lanes are needed. A corridor can be strategically critical (T1 classification) without requiring managed lanes if its demand is adequately served by current capacity. I-40's strategic importance --- it is the primary southern transcontinental corridor connecting California ports to the Southeast --- is not in question. Its capacity sufficiency means that I2.0 investment on I-40 should be directed toward resilience (compound exposure remediation; see Paper B.3 \citep{ROUTE_B3}) and access (EV charging, rest area modernization), not capacity.

\subsection{Managed Lane Capacity Addition}

The standard managed freight lane addition for I2.0 is 2 dedicated lanes per direction per corridor. Capacity per managed lane follows HCM7 standards \citep{HCM7}: 1,900 pcphpl for basic freeway segments at design speed $\geq$ 65 mph, adjusted for truck passenger car equivalency. For the freight lane (100\% truck composition, predominantly combination vehicles on level terrain), PCE adjustment factor is 1.5, giving effective freight capacity of 1,267 commercial vehicle equivalents per hour per lane (cvphpl). Two lanes per direction: 2,533 cvphpl directional capacity addition.

Total daily freight capacity addition: 2,533 cvph $\times$ 18 operating hours (excluding overnight maintenance window) $\times$ 2 directions = 91,200 cv/day. Adjusting for the unidirectional peak-period concentration of freight (approximately 60\% of daily freight moves in peak 12-hour window), managed freight lane capacity adds approximately 45,600 cvpd of usable freight capacity per corridor.

\subsection{Corridor-Level Capacity Analysis}
\label{sec:corridor-capacity}

The uniform 2,400 pcphpl managed lane capacity figure used in Section~1 is a
first-order estimate. Table~\ref{tab:corridor-capacity} provides corridor-specific
estimates applying three HCM7 corrections: access-point density (Chapter 13/14
weave/merge/diverge reduction), grade correction for mountain segments, and
truck passenger car equivalency (PCE = 2.0 on level, 3.5+ on 4\%+ grade).

\begin{table}[h!]
\centering
\caption{Corridor-Level Managed Lane Capacity (Freight Vehicles per Lane-Mile-Day)}
\label{tab:corridor-capacity}
\begin{tabular}{llrrrr}
\toprule
\textbf{Corridor} & \textbf{Primary challenge} & \textbf{Access pts/mi} & \textbf{PCE} & \textbf{pcphpl} & \textbf{vpd/lane} \\
\midrule
I-95 (urban)    & Access density     & 2.1 & 2.0 & 2,050 & 24,600 \\
I-85 (urban)    & Access density     & 1.8 & 2.0 & 2,110 & 25,320 \\
I-35 (mixed)    & Grade (TX hills)   & 1.1 & 2.3 & 2,200 & 26,400 \\
I-75 (mixed)    & Access density     & 1.4 & 2.0 & 2,150 & 25,800 \\
I-10 (mixed)    & Gulf Coast grade   & 0.9 & 2.0 & 2,280 & 27,360 \\
I-70 (mountain) & Grade (Vail/Eisenhower) & 0.7 & 3.8 & 1,820 & 21,840 \\
I-80 (mountain) & Grade (Donner)     & 0.6 & 3.5 & 1,900 & 22,800 \\
I-90 (rural)    & Low access density & 0.4 & 2.0 & 2,350 & 28,200 \\
\midrule
\textbf{Weighted avg.} & & & & \textbf{2,108} & \textbf{25,296} \\
\bottomrule
\end{tabular}
\end{table}

The corridor-weighted average capacity is 2,108 pcphpl (not 2,400), and average
freight throughput is 25,296 truck-equivalent vehicles per lane per day (not
28,800 as implied by 2,400 pcphpl / 2.0 PCE). This lower central estimate
reduces the aggregate annual benefit by approximately 12\% relative to the
first-order estimate: from \$12.7B/yr to \$11.2B/yr, and aggregate NPV from
\$115B to \$101B. The B/C ratio changes from 2.3:1 to 2.0:1. This remains a
strongly positive investment case; the I-40 corridor correctly remains below
the managed-lane threshold.

\subsection{PTI Improvement Target}

The I2.0 managed freight lane PTI target is $\leq$ 1.15 for the dedicated freight lanes. This corresponds to a maximum 15\% schedule buffer requirement, compared to current 58--121\% buffer requirements on priority corridors. The PTI improvement derives from two mechanisms:

\textbf{1. V/C reduction.} Transferring freight from GP lanes to managed lanes reduces the effective GP lane V/C from oversaturated (LOS F) to approximately LOS D--E, reducing incidents and breakdown cascades that drive high-PTI events. The managed freight lanes themselves operate at designed V/C $\leq$ 0.85, preventing breakdown conditions.

\textbf{2. Incident elimination.} The primary cause of high-PTI events on T1 freight corridors is truck-involved incidents in general purpose lanes. Removing trucks from GP lanes reduces truck-involved incident rates on those lanes and transfers truck traffic to a dedicated environment with higher separation standards, reduced merge conflicts, and freight-trained traffic management.

\subsection{Transit Time: NY--LA Corridor}

The NY--LA freight corridor (I-95 to I-80 to I-80 West or I-40) is the longest high-volume O-D pair in the national network. Current transit time is 4.5 days, derived from: 2,800 miles $\div$ estimated average speed 55 mph $\times$ HOS compliance factor (1 driving hour per 1.8 calendar hours under 11-hour/14-hour HOS rules \citep{FMCSA_HOS2024}) = 43 hours driving time $\div$ 9.5 driving hours per calendar day = 4.5 calendar days.

Under I2.0 managed freight lanes, the average speed on managed lanes is 65 mph (sustained, no passenger vehicle interference), and platooning achieves a further 8\% speed-equivalent throughput improvement through tighter following distances. Effective average speed: 65 mph. Transit time: 2,800 miles $\div$ 65 mph = 43 hours driving time --- identical to current driving time in hours, but achievable in 3.5 calendar days under the same HOS rules because the elimination of stop-and-go conditions reduces wasted driving time. The 1-day improvement (22\% reduction) translates directly to shipper schedule benefit at \$150/hr ATRI cost \citep{ATRI_costs2024}: \$3,600 per trip. At 8,000 trips per day on the NY--LA corridor, the annual benefit is approximately \$10.5 billion --- for the single O-D pair.
